Seja \(f\colon \mathbb{R}\to \mathbb{R}\) uma função. Suponha que para todo \(\varepsilon >0\), existe uma função \(g\colon \mathbb{R}\to (0,\infty)\) tal que para cada par \((x,y)\) de números reais, se \(|x-y|<\text{min}\{g(x),g(y)\}\), então \(|f(x)-f(y)|<\varepsilon\). Prove que \(f\) é o limite pontual de uma sequência de funções contínuas \(h_n\colon \mathbb{R}\to \mathbb{R}\), ou seja, existe uma sequência \(h_1,h_2,\dots,\) de funções contínuas \(h_n\colon \mathbb{R}\to \mathbb{R}\) tal que \(\lim_{n\to \infty}h_n(x)=f(x)\) para todo \(x\in \mathbb{R}\).
